% ===== latex-recomposed:document =====
% flm: style file flmrecomposed.32f07cef.sty
\begin{lemma}\flmLabels{\label{ref1}}\noproofref A lemma.
\end{lemma}

\begin{exercise}\flmLabels{\label{ref2}}
An exercise.
\end{exercise}

\flmRef[abbreviated]{ref1}, \flmRef[lowercase]{ref1}, \flmRef[capital]{ref1}, \flmRef[S]{ref1} and \flmRef[short]{ref2}.%%
% ===== latex-recomposed:style flmrecomposed.32f07cef.sty =====
\NeedsTeXFormat{LaTeX2e}
\ProvidesPackage{flmrecomposed.32f07cef}
% --- packages ---
% flm.latex.pkg.phfthm
\RequirePackage[thmset=empty]{phfthm}
% flm.latex.pkg.zref-clever
\RequirePackage{zref-clever}
% --- commands and environments the library defines for the document ---
% flm.latex.def.theorems
% Theorems. \noproofref, written after a statement's labels, says that no proof
% of the statement is coming; defined here as nothing where no package (phfthm)
% defined it first.
\ProvideDocumentCommand\noproofref{}{}
% The kinds of a class are set in the theorem style of the class's name. The
% shipped class definitionlike is set in amsthm's definition style (upright);
% declare a theorem style of that name before loading this file to change it.
\ifcsname th@definitionlike\endcsname\else\newtheoremstyle{definitionlike}{}{}{\normalfont}{}{\bfseries}{.}{5pt plus 1pt minus 1pt}{}\fi
% flm.latex.def.theorems.authored
% \flmProofTitle{<title>}: what a proof's heading sets for a title the document
% wrote, which is written as its argument inside the proof environment's
% optional argument. phfthm reads that argument as the thing the proof proves
% and words the heading itself (Proof of <title>), so by default the title
% alone. Redefine \flmProofTitle to set a proof's title differently.
\ProvideDocumentCommand\flmProofTitle{m}{#1}
% flm.latex.def.refs
% \flmRef[<options>]{<label>}: a reference that gives the name of its target
% ("Section 2"), set by \zcref of the zref-clever package.
%
% The options are the ones the document wrote, in FLM's own vocabulary. This
% definition translates each of them by its key -- a value written after the
% key is ignored, as in FLM -- and drops every key it has no translation for:
%
%   noprefix    -> noname                 the number or title alone, with no
%                                         name before it ("2", not
%                                         "Section 2")
%   S           -> S                      the first name capitalized and
%                                         unabbreviated, for a reference that
%                                         opens a sentence
%   capital     -> cap=true,abbrev=false  the names capitalized ("Lemma 1")
%   lowercase   -> cap=false,abbrev=false the names in lowercase ("lemma 1")
%   abbreviated -> abbrev=true            the abbreviated names ("Lem. 1")
%                                         where the reference type declares
%                                         them, the full names otherwise
%   anything else -> dropped              the default wording; a variant of
%                                         your own naming has no counterpart
%                                         in zref-clever, which would stop on
%                                         a key it does not know
%
% Redefine \flmRef to set references differently -- another package, another
% wording, your own variants, no link. Your definition sees the options as the
% document wrote them, in FLM's vocabulary, and maps them as it likes.
%
% The translation uses expl3, LaTeX's programming layer, which is part of the
% format and needs no package; its names read only between \ExplSyntaxOn and
% \ExplSyntaxOff. \keyval_parse:NNn reads the option list as LaTeX reads any
% key-value list: it hands a key written alone to the first macro below and a
% key with its value to the second, which drops the value. \str_case:nnF looks
% the key up in the table and appends what \zcref is to be given, or nothing.
% The macros with @ in their names are internals, not a customization point:
% redefine \flmRef, not its parts.
\ExplSyntaxOn
\def\flm@ref@add#1{\tl_put_right:Nn \flm@ref@options {#1,}}
\def\flm@ref@key#1{\str_case:nnF {#1}
  {{noprefix}{\flm@ref@add{noname}}
   {S}{\flm@ref@add{S}}
   {capital}{\flm@ref@add{cap=true,abbrev=false}}
   {lowercase}{\flm@ref@add{cap=false,abbrev=false}}
   {abbreviated}{\flm@ref@add{abbrev=true}}}
  {}}
\def\flm@ref@pair#1#2{\flm@ref@key{#1}}
\def\flm@ref@zcref#1#2{\zcref[#1]{#2}}
\ProvideDocumentCommand\flmRef{O{} m}{%
  \def\flm@ref@options{}%
  \keyval_parse:NNn \flm@ref@key \flm@ref@pair {#1}%
  \exp_args:No \flm@ref@zcref {\flm@ref@options} {#2}}
\ExplSyntaxOff
% flm.latex.def.labels
% \flmLabels{<\label commands>}: the labels of one place in the document.
% In text it runs them as written. On a line of a display-math
% environment, where amsmath keeps the last label of a line and loses the
% ones before it, it runs the first and records every further one as an alias:
% a line per alias is written to the .aux file after the last page and read on
% the next run, making the alias refer to what the first label refers to
% (for \ref and for \zcref alike). Redefine \flmLabels to change what a
% group of labels does; \def\flmLabels#1{#1} restores plain LaTeX behavior,
% amsmath's "Multiple \label's" message included.
\ProvideDocumentCommand\flmLabels{m}{%
  \ifmmode\expandafter\flm@labels@math\else\expandafter\@firstofone\fi{#1}}
% Reading the group: the first \label as written, the rest as aliases of it.
\def\flm@labels@math#1{\def\flm@lbl@first{}\flm@lbl@scan#1\flm@lbl@stop}
\def\flm@lbl@scan{\@ifnextchar\label\flm@lbl@take\flm@lbl@rest}
\def\flm@lbl@take\label#1{%
  \ifx\flm@lbl@first\@empty
    \def\flm@lbl@first{#1}%
    \label{#1}%
  \else
    \expandafter\flm@recordalias\expandafter{\flm@lbl@first}{#1}%
  \fi
  \flm@lbl@scan}
\def\flm@lbl@rest#1\flm@lbl@stop{#1}
% Collecting the aliases. A name is recorded once however often the line is
% typeset (amsmath measures a line before setting it).
\def\flm@aliaslist{}
\def\flm@recordalias#1#2{%
  \ifcsname flm@alias@#2\endcsname\else
    \expandafter\gdef\csname flm@alias@#2\endcsname{}%
    \g@addto@macro\flm@aliaslist{\flm@alias@entry{#2}{#1}}%
  \fi}
% Writing them, from the hook that runs after the last page has been shipped
% out: every \newlabel of that page has reached the .aux file by then, so the
% alias lines come after them. \AtEndDocument would run too early.
\def\flm@alias@write#1#2{\immediate\write\@auxout{\string\flm@aliaslabel{#1}{#2}}}
\AddToHook{enddocument/afterlastpage}{\let\flm@alias@entry\flm@alias@write\flm@aliaslist}
% Reading them back: the alias refers to what the original refers to. The
% assignments are global because the .aux file is read inside a group.
\def\flm@aliaslabel#1#2{%
  \ifcsname r@#2\endcsname
    \global\expandafter\let\csname r@#1\expandafter\endcsname\csname r@#2\endcsname
  \fi
  \ifcsname Z@R@#2\endcsname
    \global\expandafter\let\csname Z@R@#1\expandafter\endcsname\csname Z@R@#2\endcsname
  \fi}
% --- definitions written from the configuration of FLM features ---
% th@theoremlike (the theorems feature's, written from its configuration): declared unless \th@theoremlike is already defined, so a definition made earlier under that name is kept
\ifcsname th@theoremlike\endcsname\else\newtheoremstyle{theoremlike}{}{}{\itshape}{}{\bfseries}{.}{5pt plus 1pt minus 1pt}{}\fi
% lemma (the theorems feature's, written from its configuration): declared unless \lemma is already defined, so a definition made earlier under that name is kept
\ifcsname lemma\endcsname\else\phfMakeTheorem[thmstyle=theoremlike,counter=phfthmcounter]{lemma}{Lemma}\zcRefTypeSetup{lemma}{Name-sg={Lemma},name-sg={lemma},Name-pl={Lemmas},name-pl={lemmas},Name-sg-ab={Lem.},name-sg-ab={Lem.},Name-pl-ab={Lems.},name-pl-ab={Lems.}}\AddToHook{env/lemma/begin}{\zcsetup{countertype={lemma=lemma}}}\fi
% th@definitionlike (the theorems feature's, written from its configuration): declared unless \th@definitionlike is already defined, so a definition made earlier under that name is kept
\ifcsname th@definitionlike\endcsname\else\newtheoremstyle{definitionlike}{}{}{\itshape}{}{\bfseries}{.}{5pt plus 1pt minus 1pt}{}\fi
% exercise (the theorems feature's, written from its configuration): declared unless \exercise is already defined, so a definition made earlier under that name is kept
\ifcsname exercise\endcsname\else\phfMakeTheorem[thmstyle=definitionlike,counter=phfthmcounter,proofref=false]{exercise}{Exercise}\zcRefTypeSetup{exercise}{Name-sg={Exercise},name-sg={Exercise},Name-pl={Exercises},name-pl={Exercises}}\AddToHook{env/exercise/begin}{\zcsetup{countertype={exercise=exercise}}}\fi
\endinput
